%% =====================================================================
\section{Discussion}
\label{sec:discussion}
%% =====================================================================

\subsection{Why a Layered Defense Is Necessary}
No single control is sufficient. Cache locking is the only mechanism that addresses the
root cause---it restores the $N{:}1$ query ratio to $1{:}1$, as the coalescing ratio of
4.7 shows---but it operates below the proxy and cannot relieve congestion of the
connection path. Stale-while-revalidate makes locking \emph{fast} by giving losers an
immediate answer, but only if a stale copy exists, which is why pre-warming matters.
Application-layer rate limiting sheds database work but acts too late: by the time the
application rejects a request, the proxy and the application have already paid for
it. Edge rate limiting protects the proxy but, on its own, would not fix the
redundant-query amplification of whatever traffic it admits. The correct design
composes them in the order requests encounter them: the edge limiter caps the flood,
cache locking coalesces the surviving misses, and stale-while-revalidate serves the
coalesced waiters without delay. This mirrors the prevention--detection--response
framing urged for DDoS defense \cite{zargar2013} and the overload-control principle of
shedding load as early as possible \cite{welsh2001,beyer2016}.

\subsection{Detecting Harm versus Detecting Novelty}
The detection results argue against treating the two detectors as substitutes, and
they suggest a sharper way to describe the difference than ``rules versus machine
learning''. The rule is a detector of \emph{harm}: it fires only when the database is
demonstrably overloaded, so it is silent on absorbed bursts, silent on defended
attacks, and silent when idle. The IsolationForest is a detector of \emph{novelty}: it
fires whenever the traffic vector leaves the region it was trained on, which happens
during attacks, but also during idleness and partly during defended attacks. Both
questions are useful. Harm is what an on-call engineer must act on; novelty is what a
security analyst wants to know about even when the defenses are holding, because a
defended attack is still an attack. A practical deployment should therefore route the
two signals differently---the rule to paging, the forest to a security dashboard---and
should give the forest an explicit notion of ``no traffic'' (for example, by not
scoring intervals with fewer than a minimum number of requests) to avoid the
closed-world false alarms we observed \cite{sommer2010}.

The analysis also shows that a plausible-sounding signal can be misleading. The
stampede narrative suggests that the cache-miss ratio should spike during an attack,
and it was included in the rule for that reason. In this testbed it does not: the short
TTL and uniform catalog make the normal miss ratio high (0.62), and during the attack
the high-volume, repeatedly requested hot keys are hits most of the time between
expiries, so the ratio actually falls (0.51). Equation~\eqref{eq:baseline} predicts the
first effect from first principles. The lesson for practitioners is to validate every
detection feature against a model or a measurement of normal behavior, and to prefer
signals that directly measure the protected resource (database attempts, pool
exhaustion) over proxies for it.

\subsection{Engineering Lessons from Building the Testbed}
\label{sec:lessons}
Several of the most instructive results emerged as bugs during development. Table
\ref{tab:bugs} summarizes them; each is a reusable lesson for anyone building or
defending a cached service.

\begin{table}[t]
\centering
\caption{Defects found during development, their causes, and the general lesson.}
\label{tab:bugs}
\footnotesize
\begin{tabularx}{\linewidth}{@{}L{2.9cm}L{3.3cm}L{3.3cm}Y@{}}
\toprule
Symptom & Root cause & Fix & General lesson \\
\midrule
Defended shop still slow; attacker throttled legitimate users & Limiter keyed on the proxy's address; all clients shared one budget & Proxy overwrites \texttt{X-Forwarded-For} with the real peer address & Per-client limits behind a proxy need trusted client identity \\
Defended latency up to ${\approx}2.5$~s & Lock losers queued behind the winner & Serve the stale copy immediately (SWR); pre-warm stale cache & Coalescing must not turn into queueing \\
Shop slow (300--700~ms) even with DB protected & Flood still congested Nginx and the app's threads & Edge \texttt{limit\_req}, toggled live via include + reload & Protect the request path, not only the database \\
Single false alarm at start of baseline & Cold cache: every request misses, looks like a stampede & Suppress alerts for 3~s warm-up; train ML only after it & Distinguish start-up transients from attacks \\
ML stopped recognizing attacks & Cold-start spike polluted the ``normal'' training set & Skip warm-up; standardize features & Training data hygiene determines unsupervised accuracy \\
Occasional one-interval ML alert & Noise in a single sample & Two-interval debounce; contamination 0.03$\rightarrow$0.01 & Debounce novelty detectors \\
Metric totals split and inconsistent & Several Gunicorn workers, each with its own counters & One worker, 200 threads & In-process metrics need a single process (or multiprocess mode) \\
Shop stuck ``red'' after attack ended & Browser fetches without timeout piled up & 4~s \texttt{AbortController} timeout; non-overlapping refresh & Clients need timeouts too, or they amplify outages \\
Dashboard felt laggy & 1~Hz sampling & 2~Hz sampling and streaming & Rates need enough samples per window; 2~Hz suffices \\
\bottomrule
\end{tabularx}
\end{table}

\subsection{Deployment Guidance}
For operators of cached, containerized web services, the results translate into the
following checklist:
\begin{enumerate}[leftmargin=*]
\item \textbf{Coalesce regeneration of hot keys} with a single-flight lock or lease,
and release locks with a token check \cite{kleppmann2016}.
\item \textbf{Keep a stale copy} of hot entries and serve it to lock losers
(stale-while-revalidate); pre-warm it at start-up so the defense works from the first
miss.
\item \textbf{Rate-limit at the edge}, not only in the application, and return 429
rather than 503 for intentional shedding (\texttt{limit\_req\_status 429}).
\item \textbf{Forward trusted client identity} from the proxy (overwrite, never
append, \texttt{X-Forwarded-For} at the outermost trusted hop) so per-client limits are
meaningful.
\item \textbf{Choose pool behavior deliberately.} A fail-fast pool converts overload
into immediate errors (good for shedding, bad for goodput under stampede); a bounded
wait with a short timeout may preserve more useful work.
\item \textbf{Monitor the protected resource directly}: database attempts, pool
exhaustion, and failures, together with latency. Validate any proxy signal (such as
miss ratio) against normal behavior before alerting on it.
\item \textbf{Run both a harm detector and a novelty detector}, route them to
different audiences, and train the novelty detector on a clean, sufficiently long
window that excludes start-up and idle periods.
\item \textbf{Treat TTLs as public.} Hit/miss timing reveals them; do not rely on
obscurity or jitter for protection.
\end{enumerate}
All of these are realizable with standard components---Redis, Nginx, and a metrics
pipeline---and none requires changes to application semantics beyond accepting bounded
staleness.

\subsection{Security of the Testbed and Ethical Considerations}
\label{sec:ethics}
CacheCollapse is a tool for generating denial-of-service traffic, and it was designed
to be safe by construction. The attack generator targets only the stack's own proxy on
a private network; no external system is ever contacted; and all data are synthetic.
The control API that launches attacks is intentionally unauthenticated for classroom
use, and Grafana ships with default credentials and anonymous viewing. These choices
are acceptable on a laptop but not on any reachable host. The project's deployment
plan accordingly restricts exposure to known addresses via a cloud firewall, requires
changing default passwords, and recommends local demonstration as the safest option;
a helper script that publishes the dashboard through a temporary tunnel should be used
only for short, supervised sessions, since anyone with the link can launch the attack
against the demo itself. The techniques described here are well known to defenders and
attackers alike; we believe that publishing a reproducible, defensive-first account of
them, with mitigations and detection, benefits defenders more than it helps attackers.

\subsection{Threats to Validity}
\label{sec:validity}
\paragraph{Construct validity (are we measuring what we claim?).} The log analysis
exposed several measurement subtleties that affect interpretation:
\begin{itemize}[leftmargin=*]
\item \emph{Database ``queries'' are attempts.} The counter is incremented before a
pool connection is obtained, so under overload it includes attempts rejected with 503.
We therefore report ``attempts'' and derive completed queries separately
(Section~\ref{sec:mechanics}).
\item \emph{Failures are logged as per-interval counts}, not per-second rates. With
0.5~s intervals the two differ by a factor of two; all rates in this paper divide by
the measured interval.
\item \emph{Hits and misses overlap under defense}, because a lock loser served stale
data is counted as both a miss and a hit; the defended miss ratio is therefore biased
low, while the coalescing ratio is unaffected.
\item \emph{Latency is measured inside the application.} It excludes time spent in the
proxy and requests rejected at the edge; the storefront and the test suite provide the
complementary end-to-end view.
\item \emph{p95 latency is quantized} to histogram bucket bounds and is therefore
reported only as supporting evidence.
\item \emph{Edge rejections are invisible to the application metrics} and, with Nginx
defaults, are HTTP~503 at the client.
\item \emph{Per-container CPU and memory} were part of the original specification but
are not collected by the final monitor (Prometheus scrapes only the application), so
no resource-utilization results are reported.
\item \emph{Defense configuration is not logged per interval.} The mode ``protected''
means ``at least one defense on''; the ablation therefore relies partly on observation
and on the session history rather than on per-interval labels.
\end{itemize}

\paragraph{Internal validity.} The attacker, the victim, and the monitor share one
host's CPU, so heavy attack traffic also slows the instrumentation; the monitor's
sampling interval stretches under load (median 0.53~s overall, longer during attacks).
Mode labels can lag traffic changes by up to one interval, which affects detection
latencies at the 0.5~s scale but not episode-level conclusions.

\paragraph{External validity.} The testbed is a single-host Docker Desktop deployment
whose port forwarding adds latency under heavy connection load, so absolute latencies
are environment-specific while relative effects are robust. Traffic is synthetic:
popularity is uniform rather than Zipf-like \cite{breslau1999}, database cost is
modelled with \texttt{pg\_sleep}, and both normal and attack traffic originate from a
single load-generator address. The last point has two consequences: the edge limiter
also throttles the synthetic normal traffic (which is why defended database load falls
below the baseline), and per-address limiting is easier than it would be against a
distributed attacker using many addresses. Legitimate-user availability under defense
is therefore evidenced by the storefront and test case~16, which originate from a
different address, rather than by the synthetic background traffic.

\paragraph{Conclusion validity.} The main results come from one final session.
Consecutive intervals are autocorrelated, so bootstrap intervals computed over
intervals understate uncertainty; the episode-level and cross-session results
(Figures~\ref{fig:dev} and~\ref{fig:detlat}) provide the more conservative evidence and
are consistent with the interval-level estimates.

\subsection{Limitations}
\label{sec:limitations}
Beyond the validity threats above, the study has the following limitations. The
IsolationForest requires a clean baseline window, so a system attacked from start-up
trains on contaminated data or not at all (as in the first episode of the final
session); the always-available rule mitigates but does not eliminate this. Thresholds
($K$, $\tau_L$, $\tau_m$, limiter rates) were tuned on the testbed and need site-specific
calibration. The edge limiter is per source address and would need to be complemented
by aggregate or behavioral limiting against a widely distributed attacker. Finally, the
study evaluates a single cache, a single database, and a single application replica;
interactions with replication, sharding, and auto-scaling are left to future work.
